\subsection{2.1 主动换道能力}

主动换道能力占 80\%，是高速领航辅助最核心的部分。它下面有三个二级指标：

\subsubsection{2.1.1 车道占用换道}

这一块考的是：本车道前方不适合继续走了，系统能不能识别出来，并选择更合理的处理方式。它又分三类。

\begin{enumerate}
  \item \textbf{前方遇静止障碍物}：比如弯道出口有静止乘用车，或者直道虚线车道内有静止障碍物。测试车速有 60、80、100 km/h 等。系统如果能够安全换道避让，效率得分更高；如果只是在本车道刹停，也能体现安全性，但效率分会低。发生碰撞、横向加速度过大，基本就是 0 分。
  \item \textbf{前方遇缓行车辆}：比如弯道中前方乘用车以 40 km/h 慢行，或仿真场景里的缓行卡车、拥堵车流旁边车道变顺畅。这里不是只看不撞，还看系统能不能判断“继续跟着太低效”，并在条件允许时主动换道超越。若本车道跟车，安全上可以得分，但效率不如换道。
  \item \textbf{前方遇切入车辆}：相邻车道车辆突然切入本车道，测试车速 80 km/h，目标车 70 km/h，切入时距 1.5 s。这里主要看系统对切入趋势的预测能力，不能等目标完全压进本车道才反应。
\end{enumerate}

这一组评分共同看三件事：\textbf{安全、舒适、效率}。一般安全权重 40\%，舒适 30\%，效率 30\%。舒适主要看最大纵向减速度和最大横向加速度： 

\begin{itemize}
  \item 最大纵向减速度 ≤ 2 m/s² 且最大横向加速度 ≤ 2 m/s²，舒适性最好； 
  \item 纵向减速度到 2-4 m/s²，或横向加速度到 2-3 m/s²，会降舒适分； 
  \item 横向加速度 > 3 m/s²、碰撞、压线等，会直接导致很差结果。
\end{itemize}

假设测试场景是：高速领航激活，前方车道有静止车，系统要处理。

\begin{itemize}
  \item 如果车辆\textbf{成功换道避开障碍物}，而且最大纵向减速度 ≤ 2 m/s²、最大横向加速度 ≤ 2 m/s²，那么：
\end{itemize}

\texttt{安全 100 × 40\% + 舒适 100 × 30\% + 效率 100 × 30\% = 100 分}

\begin{itemize}
  \item 如果也成功换道了，但动作比较猛，比如最大纵向减速度到了 3 m/s²，或者横向加速度到了 2.5 m/s²，那么舒适性降成 60 分：
\end{itemize}

\texttt{安全 100 × 40\% + 舒适 60 × 30\% + 效率 100 × 30\% = 88 分}

\begin{itemize}
  \item 如果系统没有换道，只是在本车道刹停，而且刹得比较平顺：
\end{itemize}

\texttt{安全 100 × 40\% + 舒适 100 × 30\% + 效率 0 × 30\% = 70 分}

\subsubsection{2.1.2 突发情况换道}

这一块考的是高速 NOA 面对非正常道路事件的处理能力。包括道路施工和事故区域。

\begin{itemize}
  \item \textbf{前方遇道路施工}：典型是直道上有 5 个斜置锥桶，车速 80 km/h。系统要识别施工区域并完成换道。如果相邻车道有障碍车，规范反而不鼓励硬换道，这时安全刹停也可以得分。这个设计很重要，说明 C-ICAP 不是盲目奖励“敢变道”，而是奖励“该变才变，不能变就安全停”。
  \item \textbf{前方遇事故区域}：包括横置/斜置事故车、斜置货车侵占车道、中雨条件下前方追尾车辆并出现打伞行人等。这里考的比普通避障更难，因为目标形态不规整，可能还有行人、雨天低附着、相邻车道慢车流。系统既要识别事故风险，又要判断能不能绕行，不能为了通过而制造更大风险。
\begin{enumerate}
  \item \textbf{普通施工/事故场景}多数突发情况换道场景，沿用“前方遇静止障碍物”的打分逻辑，也就是看三项：
\end{enumerate}

典型判定是：

最后用：

\texttt{单项得分 = 安全 × 40\% + 舒适 × 30\% + 效率 × 30\%}

\begin{enumerate}
  \item \textbf{特殊施工场景：相邻车道有障碍车}这个场景不看效率，只看安全和舒适，因为旁边有障碍车，不鼓励系统硬变道。
\end{enumerate}

判定表是：

\subsubsection{2.1.3 导航路线换道}

这一块是高速 NOA 区别于普通 LCC/ACC 的关键。因为导航路线要求它上下匝道。

\begin{itemize}
  \item \textbf{主路汇入匝道}：包括无干扰车辆、匝道口有静止车辆、最右侧车道有连续行驶车流。测试车速覆盖 80、100、120 km/h。它考的是系统能不能提前规划车道位置，选择正确匝道，并在匝道口有障碍时安全处理。
  \item \textbf{匝道汇入主路}：仿真场景中，主路左侧邻车道后方有非连续行驶车流，车流速度 70 km/h、间距 100 m。这里重点看系统能不能找到安全间隙并入主路。评分里有一个“安全时间域”概念：与前车和后车的安全时间域都 ≥ 0.5 s，安全分更高；只要其中一侧 < 0.5 s，安全得分会降。
\end{itemize}

导航路线换道下面主要看两类：

\begin{enumerate}
  \item \textbf{主路汇入匝道怎么打分}主路汇入匝道又分三种场景：
\end{enumerate}

无干扰车辆时，系统要从主路按导航进入目标匝道。每个速度点重复 3 次，取表现最差的一次。速度点是：

评分分三块：\textbf{安全 40\% + 舒适 30\% + 效率 30\%}。

如果系统成功汇入匝道，并且最大纵向减速度 ≤ 2 m/s²、最大横向加速度 ≤ 2 m/s²，就是最理想状态，安全、舒适、效率都 100 分。

如果也能汇入匝道，但动作比较急，比如纵向减速度在 2-4 m/s²，或者横向加速度在 2-3 m/s²，安全和效率还能拿 100，但舒适会降到 60。

如果系统提醒驾驶员介入，安全只有 50，舒适和效率都是 0。

如果驶出车道、发生碰撞、无法进入匝道缓冲区，或者最大横向加速度 > 3 m/s²，就是 0 分。

\begin{enumerate}
  \item \textbf{匝道口有静止车辆怎么打分}这个场景不是鼓励它硬进匝道，而是看系统能不能识别匝道口风险并安全处理。
\end{enumerate}

评分只看：\textbf{安全 60\% + 舒适 40\%}。

也就是说：\textbf{匝道口有静止车时，安全刹停是可以得高分的。}

\begin{enumerate}
  \item \textbf{连续车流场景怎么打分}“主路最右侧车道有连续行驶车流”是仿真加分项。测试车速 120 km/h，右侧车流 70 km/h，车间距 100 m。
\end{enumerate}

这里最关键的是“安全时间域”：本车变道进入目标车道时，与前车、后车都要有足够时间间隔。
